\documentclass[12pt,letterpaper]{article}
\usepackage[utf8]{inputenc}
\usepackage[spanish]{babel}
\usepackage[margin=2.5cm]{geometry}
\usepackage{booktabs}
\usepackage{longtable}
\usepackage{array}
\usepackage{colortbl}
\usepackage{xcolor}
\usepackage{graphicx}
\usepackage{caption}
\usepackage{enumitem}
\usepackage{titlesec}
\usepackage{parskip}
\usepackage{hyperref}
\usepackage{tabularx}
\usepackage{ragged2e}
\usepackage{graphicx}
\usepackage{float}      % <-- agregar aquí


% ── Colores corporativos de las tablas ──────────────────────────────────────
\definecolor{headerblue}{RGB}{31,73,125}
\definecolor{rowlightblue}{RGB}{217,235,255}
\definecolor{headerorange}{RGB}{197,90,17}
\definecolor{rowlightorange}{RGB}{252,228,214}
\definecolor{darkgray}{RGB}{64,64,64}

% ── Formato de títulos ───────────────────────────────────────────────────────
\titleformat{\section}{\large\bfseries}{}{0em}{}
\titleformat{\subsection}{\normalsize\bfseries}{}{0em}{}

\begin{document}
\begin{titlepage}
    \centering
    \vspace*{2cm}
    
    {\large Universidad de los Andes \par}
    \vspace{0.3cm}
    {\large Ingeniería de Sistemas y Computación \par}
    \vspace{0.3cm}
    {\large Architecting Digital Enterprises \par}
    
    \vspace{2.5cm}
    
    {\LARGE \textbf{Caso de Estudio: Todo Tarjetas S.A.} \par}
    \vspace{0.8cm}
    {\Large Taller 2 --- Cultura, Validación UX y Arquitectura de Solución \par}
    
    \vspace{3cm}
    
    {\large
        Esteban Castelblanco \\[0.3cm]
        Julián Baquero \\[0.3cm]
        Nicolás Casas \par
    }
    
    \vfill
    
    {\large Bogotá, Colombia \par}
    \vspace{0.3cm}
    {\large 2026 \par}
    
    \vspace{2cm}
\end{titlepage}


% ============================================================
\section*{1. Definición de Cultura para Digital}
% ============================================================

% ------------------------------------------------------------
\subsection*{a)  Defina y justifique el modelo estructural del nuevo negocio (SpinOff Vs Nueva Línea de Negocio)  }
% ------------------------------------------------------------

Se propone una \textbf{Nueva Línea de Negocio Digital} dentro de Todo Tarjetas S.A.\ como
modelo para escalar los dos CVPs definidos. La transformación se centra en dos pilares:

\begin{itemize}
  \item \textbf{CVP 1:} un asistente GenAI omnicanal que elimina la fricción del call center.
  \item \textbf{CVP 2:} un proceso de originación 100\,\% digital con tarjeta virtual instantánea.
\end{itemize}

Ambos CVPs se benefician de los activos existentes de Todo Tarjetas. Se propone una Nueva
Línea de Negocio Digital para habilitar los dos CVP priorizados, aprovechando los activos
existentes de la compañía.

\subsubsection*{Aprovechamiento de Activos Existentes de Todo Tarjetas S.A.}

\begin{itemize}
  \item \textbf{Infraestructura de procesamiento disponible:} el CVP1 se implementa sobre el
    switch de procesamiento actual, que ya opera 2\,500 prospectos semanales, eliminando la
    necesidad de nueva infraestructura tecnológica.
  \item \textbf{Motor de riesgo reutilizable:} el CVP2 toma como base el motor de reglas
    existente y le incorpora una capa de ML, evitando el desarrollo de un sistema de evaluación
    crediticia desde cero.
  \item \textbf{Base de clientes consolidada:} las relaciones vigentes con bancos emisores y
    comercios afiliados permiten la adopción inmediata de ambos CVP sin incurrir en costos
    adicionales de adquisición de clientes.
  \item \textbf{Licencias regulatorias vigentes:} Todo Tarjetas opera actualmente bajo
    supervisión de la SFC. Constituir un Spin-Off implicaría tramitar una nueva licencia,
    proceso que puede tomar entre 12 y 18 meses y genera costos adicionales significativos
    antes de poder iniciar operaciones.
\end{itemize}

\subsubsection*{Ventajas sobre el SpinOff — Análisis Específico al Business Case}

\begin{itemize}
  \item \textbf{Velocidad de llegada al mercado:} una Nueva Línea de Negocio puede estar
    operativa en 8 a 12 meses, mientras que un Spin-Off requiere entre 18 y 24 meses. En un
    mercado de procesadoras digitales en rápida evolución, ese tiempo de diferencia representa
    una ventana competitiva que no puede perderse.
  \item \textbf{Inversión compartida:} el CAPEX de USD 530\,K proyectado en el Business Case
    está diseñado sobre una infraestructura compartida con la operación actual de Todo
    Tarjetas. Bajo un esquema de Spin-Off, esta estructura de costos no sería viable y la
    inversión requerida sería significativamente mayor.
  \item \textbf{Despliegue inmediato:} el CVP1 puede implementarse sobre los canales digitales
    existentes (App y Web) desde el primer ciclo de desarrollo. Un Spin-Off tendría que
    construir toda la experiencia de usuario desde cero, lo que implica tiempo y costos
    adicionales no contemplados.
  \item \textbf{Valor financiero en riesgo:} el VPN de USD 12,2\,M proyectado a 5 años asume
    una adopción del 60\,\% de los clientes en el primer año. Ese nivel de adopción solo es
    alcanzable contando con la base de clientes existente — un Spin-Off no tendría esa base y
    comprometería directamente el retorno financiero del proyecto.
\end{itemize}

% ============================================================
\subsection*{b)Estructure el modelo de congruencia para digital que va determinar la Cultura que se requiere para poner en funcionamiento el Digital Business Model, para esto defina conductas  observables, comportamientos y procesos para cada elemento del modelo de congruencia.}
% ── Espacio para imagen del modelo de congruencia ─────────────────────────
\begin{figure}[h!]
  \centering
  \includegraphics[width=\textwidth]{artefactos/modelo_congruencia.png}
  \caption{Modelo de Congruencia para Digital — Todo Tarjetas S.A.}
\end{figure}
\bigskip

% ============================================================

Teniendo en cuenta la premisa, se presenta el siguiente modelo de congruencia para determinar
la cultura requerida.

% ── Tabla del modelo de congruencia ────────────────────────────────────────
\begingroup
\small
\setlength{\tabcolsep}{4pt}
\renewcommand{\arraystretch}{1.3}

\begin{longtable}{%
  >{\bfseries\RaggedRight}p{1.8cm}
  >{\RaggedRight}p{2.6cm}
  >{\RaggedRight}p{2.8cm}
  >{\RaggedRight}p{2.6cm}
  >{\RaggedRight}p{2.6cm}
  >{\RaggedRight}p{2.4cm}}

\rowcolor{headerblue}
\textcolor{white}{\textbf{Elemento}} &
\textcolor{white}{\textbf{Estado Actual}} &
\textcolor{white}{\textbf{Estado Deseado}} &
\textcolor{white}{\textbf{Conductas Observables}} &
\textcolor{white}{\textbf{Comportamientos}} &
\textcolor{white}{\textbf{Procesos}} \\
\endfirsthead

\rowcolor{headerblue}
\textcolor{white}{\textbf{Elemento}} &
\textcolor{white}{\textbf{Estado Actual}} &
\textcolor{white}{\textbf{Estado Deseado}} &
\textcolor{white}{\textbf{Conductas Observables}} &
\textcolor{white}{\textbf{Comportamientos}} &
\textcolor{white}{\textbf{Procesos}} \\
\endhead

% ── Fila: Estrategia ──
\rowcolor{rowlightblue}
Estrategia &
Procesamiento transaccional tradicional. Call center como canal primario. OPEX $>$ 2\,\% del mercado. &
GenAI como palanca de autoservicio. Originación 100\,\% digital. OPEX reducido vía automatización. &
Decisiones basadas en KPIs digitales. A/B testing de prompts LLM. Priorizar reducción PQR en dashboard. &
Revisión semanal de métricas CVP. Iteración de modelos ML cada sprint. Benchmarking vs.\ fintechs. &
OKRs por CVP. Planificación ágil 90 días. Revisión de adopción mensual. \\

% ── Fila: Cultura ──
Cultura &
Aversión al riesgo operativo. Cumplimiento SLA call center. Cambios requieren múltiples aprobaciones. &
Tolerante al fallo controlado. Experimentación con GenAI en sandbox. Cultura de aprendizaje continuo. &
Compartir post-mortems sin culpa. Celebrar automatización de PQR. Proponer mejoras al bot activamente. &
Adoptar Scrum para squads CVP. Retrospectivas quincenales abiertas. Mentalidad de producto sobre proyecto. &
Post-mortems documentados. Daily de squads. Policy de no culpa por fallos en pruebas. \\

% ── Fila: Personas ──
\rowcolor{rowlightblue}
Personas &
Especialistas en call center. Operadores de scoring manual. Analistas de documentos físicos. &
Ingenieros de IA/ML y Prompt Engineering. Diseñadores UX para originación digital. Product Owners con visión de negocio. &
Rotación de agentes a squads GenAI. Upskilling en LangChain y OCR. Mentoring en Design Thinking. &
Autogestión del aprendizaje en IA. Flexibilidad ante cambios de modelo LLM. Colaboración cross-funcional. &
Programa reskilling agentes call center. Tours of duty de 6 meses. Evaluación 360° con criterio digital. \\

% ── Fila: Tareas ──
Tareas &
Resolución PQR por teléfono manual. Digitación manual de documentos. Scoring en 2–3 días hábiles. &
Autoservicio GenAI 24/7 para PQR. OCR + biometría automática KYC. Scoring automático en $<$ 60 segundos. &
Uso de LangChain Orchestrator diario. Documentación de prompts en repositorio. Feedback continuo al modelo LLM. &
Sprints de 2 semanas por CVP. CI/CD para despliegue del bot. A/B testing de flujos de originación. &
Kanban digital por squad. Definition of Done con criterio de calidad IA. Retraining del modelo cada 30 días. \\

% ── Fila: Estructura ──
\rowcolor{rowlightblue}
Estructura &
Vertical: Call Center + TI + Riesgo en silos. TI Legacy separado del negocio. Decisiones centralizadas en Gerencia. &
Squads autónomos: Squad GenAI (CVP1) + Squad Originación (CVP2). Product Owners con P\&L propio. Liderazgo distribuido. &
Decisiones del PO sin aprobación externa. Comunicación directa squads $\leftrightarrow$ core bancario. OKRs públicos por CVP. &
Equipos autónomos end-to-end. Rotación de roles entre squads. Mentoring cruzado PO $\leftrightarrow$ Ingeniero. &
Squads/Tribes por CVP. Guilds: Arquitectura · IA · UX. Chapter de Data Engineering compartido. \\

\end{longtable}
\endgroup

\bigskip

\subsubsection*{CVP1 — Asistente GenAI Omnicanal + Autoservicio}

\begin{itemize}
  \item \textbf{Estrategia:} el call center representa un costo de USD 1\,500\,000 anuales.
    El CVP1 permite recuperar USD 472\,500 al resolver automáticamente el 35\,\% de las
    solicitudes sin intervención humana. Para que esto funcione, la organización debe medir su
    éxito por la cantidad de casos resueltos digitalmente, no por el volumen de llamadas
    atendidas.
  \item \textbf{Cultura:} el asistente de inteligencia artificial cometerá errores en sus
    primeras etapas de operación. Si la organización no tiene tolerancia al fallo, cualquier
    error inicial será utilizado como justificación para cancelar el proyecto. La cultura debe
    enfocarse en la mejora continua y celebrar el aprendizaje iterativo.
  \item \textbf{Personas:} los agentes del call center son quienes mejor conocen los casos
    reales de los clientes. En lugar de sentirse desplazados, deben ser protagonistas del
    proceso: su conocimiento es insumo clave para entrenar y mejorar el modelo de inteligencia
    artificial.
  \item \textbf{Tareas:} el asistente requiere revisión y ajuste permanente. Esto implica
    equipos de trabajo que operen en ciclos cortos de dos semanas, analizando métricas de
    resolución y realizando mejoras continuas al modelo.
  \item \textbf{Estructura:} el equipo responsable del CVP1 necesita acceso directo a los
    sistemas de procesamiento y gestión de clientes sin depender de aprobaciones de comités.
    La estructura organizacional debe facilitar esa autonomía, no convertirse en un obstáculo.
\end{itemize}

\subsubsection*{CVP2 — Originación 100\,\% Digital + Tarjeta Virtual Instantánea}

\begin{itemize}
  \item \textbf{Estrategia:} actualmente se destinan USD 1\,800\,000 anuales a la producción y
    entrega física de tarjetas plásticas. La estrategia debe reorientar sus métricas de éxito:
    cada tarjeta virtual emitida representa un ahorro directo frente a cada tarjeta física
    producida.
  \item \textbf{Cultura:} los analistas de riesgo pueden percibir el scoring automático como
    una amenaza a su rol. La organización debe comunicar con claridad que el CVP2 no los
    reemplaza, sino que los libera de tareas repetitivas para que puedan concentrarse en casos
    complejos de mayor valor.
  \item \textbf{Personas:} la validación biométrica y el reconocimiento de documentos requieren
    perfiles técnicos especializados que hoy no existen dentro de la compañía. Los equipos
    deben integrar talento externo especializado con el equipo interno de riesgo, trabajando de
    manera complementaria.
  \item \textbf{Tareas:} el reconocimiento facial puede fallar en condiciones de baja
    iluminación u otros factores externos. Las responsabilidades del equipo no se limitan al
    flujo ideal de aprobación — también incluyen el control de calidad del modelo, el manejo
    de excepciones y el diseño de rutas alternas para los casos que no siguen el camino
    estándar.
  \item \textbf{Estructura:} el equipo de originación digital debe tener autonomía para ajustar
    el motor de scoring sin requerir aprobación gerencial para cada modificación. Sin esa
    autonomía, los ciclos de desarrollo se bloquean y el ritmo de mejora del CVP se detiene.
\end{itemize}

% ============================================================
\subsection*{c) Teniendo como contexto el framework de análisis de cultura digital mencionado en la figura que se presenta a continuación, y el cual fue 
revisado en clase, estructure y realice un análisis que permita definir la cultura digital requerida en Todo Tarjetas para asegurar el éxito de la 
conceptuación, construcción y adopción de los CVP que definió.     
Es importante que esta cultura habilite la transformación del negocio, movilice 
conductas observables y comportamientos, así como genere hábitos de trabajo basados en los siguientes elementos que caracterizan a las 
organizaciones maduras digitales}

% ============================================================
% ── Espacio para imagen del framework de cultura digital ────────────────────
\begin{figure}[H]
  \centering
  \includegraphics[width=\textwidth]{artefactos/framework_cultura.png}
  \caption{Framework de Cultura Digital — Rating 1 a 5 (Early, Developing, Maturing)}
\end{figure}

A continuación se analiza cada dimensión del framework de cultura digital, identificando el
estado actual, el estado objetivo, la calificación y la iniciativa clave asociada.

% ── Tabla de dimensiones de cultura digital ────────────────────────────────
\begingroup
\small
\setlength{\tabcolsep}{4pt}
\renewcommand{\arraystretch}{1.35}

\begin{longtable}{%
  >{\bfseries\RaggedRight}p{2cm}
  >{\RaggedRight}p{3.2cm}
  >{\RaggedRight}p{3.6cm}
  >{\centering}p{1.1cm}
  >{\centering}p{1.5cm}
  >{\RaggedRight}p{3.4cm}}

\rowcolor{headerblue}
\textcolor{white}{\textbf{Dimensión}} &
\textcolor{white}{\textbf{FROM}} &
\textcolor{white}{\textbf{TO}} &
\textcolor{white}{\textbf{Actual (1–5)}} &
\textcolor{white}{\textbf{Objetivo Año 2}} &
\textcolor{white}{\textbf{Iniciativa Clave}} \\
\endfirsthead

\rowcolor{headerblue}
\textcolor{white}{\textbf{Dimensión}} &
\textcolor{white}{\textbf{FROM}} &
\textcolor{white}{\textbf{TO}} &
\textcolor{white}{\textbf{Actual (1–5)}} &
\textcolor{white}{\textbf{Objetivo Año 2}} &
\textcolor{white}{\textbf{Iniciativa Clave}} \\
\endhead

% ── Agilidad ──
\rowcolor{rowlightblue}
Agilidad &
Los cambios en el sistema de procesamiento requieren entre 4 y 6 semanas de aprobación por comités de arquitectura. &
Ciclos de desarrollo de 2 semanas por CVP, despliegue continuo del asistente con control de versiones y capacidad de implementar cambios en menos de 2 días. &
2,0 & 4,0 &
Adoptar metodología Scrum. Implementar integración y despliegue continuo para el modelo de IA. Eliminar comités de aprobación para ajustes operativos. \\

% ── Apetito al riesgo ──
Apetito al riesgo &
La cultura de cumplimiento bancario genera aversión total al error. Cualquier fallo del asistente se interpreta como razón suficiente para cancelar el proyecto. &
Ambiente controlado de pruebas para el modelo de IA con tráfico limitado. Los errores se documentan como aprendizajes. Se realizan pruebas A/B de manera sistemática. &
1,5 & 4,0 &
Crear un sandbox de innovación dedicado. Establecer una política de no culpabilización por errores. Proteger un presupuesto de innovación de USD 60\,K independiente del presupuesto operativo. \\

% ── Toma de decisiones ──
\rowcolor{rowlightblue}
Toma de decisiones &
Las decisiones de evaluación crediticia las toma el Gerente de Riesgo de manera manual y sin acceso a datos en tiempo real. &
Panel de indicadores en tiempo real para el asistente y el scoring. El Product Owner toma decisiones de producto de manera autónoma. Las aprobaciones crediticias se realizan de forma automática. &
1,5 & 4,5 &
Implementar panel de indicadores por CVP en tiempo real. Capacitar a los líderes de producto en análisis de datos. Adoptar modelos de scoring explicables que generen confianza en la toma de decisiones automatizada. \\

% ── Liderazgo ──
Liderazgo &
La Gerencia de Operaciones aprueba todos los cambios en canales digitales y procesos de riesgo, centralizando las decisiones. &
Los líderes de producto tienen responsabilidad directa sobre los resultados financieros del CVP1 y CVP2. Los equipos operan con autonomía real y cuentan con acompañamiento de liderazgo. &
2,0 & 4,0 &
Implementar programa de liderazgo ágil. Otorgar autoridad real a los líderes de producto. Eliminar la capa de aprobación intermedia que ralentiza la ejecución. \\

% ── Pasión por el trabajo ──
\rowcolor{rowlightblue}
Pasión por el trabajo &
Alta rotación en el call center. Los agentes no perciben un propósito más allá de resolver el volumen diario de llamadas. &
Propósito claro: automatizar las tareas repetitivas para mejorar la experiencia del cliente. Los agentes se convierten en entrenadores del asistente y sienten orgullo por los resultados de satisfacción. &
2,5 & 4,5 &
Comunicar la misión de manera consistente: eliminar los tiempos de espera para el cliente. Involucrar activamente a los agentes del call center en el entrenamiento del modelo de inteligencia artificial. \\

% ── Estilo de trabajo ──
Estilo de trabajo &
El call center, TI y Riesgo operan de forma aislada, sin herramientas ni métricas compartidas entre áreas. &
Equipos multidisciplinarios conformados por ingenieros, analistas de riesgo y diseñadores de experiencia. Todos comparten métricas e indicadores en tiempo real. &
1,5 & 4,5 &
Implementar modelo de equipos por CVP. Estandarizar herramientas de trabajo colaborativo. Realizar sesiones mensuales de innovación con participación de todas las áreas. \\

\end{longtable}
\endgroup

% ============================================================
\subsection*{d) Costo o riesgo de no hacer la evolución cultural requerida}
% ============================================================

\subsubsection*{1. Corto Plazo (0–12 meses)}

\begin{itemize}
  \item \textbf{Pérdida de talento digital clave:} sin una cultura ágil ni ownership real sobre
    los CVPs, los perfiles especializados en modelos de lenguaje, biometría y product
    management que se contraten para el CVP1 y el CVP2 migrarán a fintechs competidoras en
    menos de 12 meses.
  \item \textbf{Adopción limitada del asistente GenAI:} si el equipo del call center no
    interioriza el cambio cultural, seguirá resolviendo manualmente las solicitudes que el CVP1
    debería atender de forma automática. El ahorro proyectado de USD 472\,500 anuales no se
    materializaría en el Año 1.
  \item \textbf{Costos operativos duplicados:} mantener el call center tradicional en paralelo
    con el asistente digital durante un período prolongado de resistencia interna generaría
    costos operativos superiores a los del modelo actual, eliminando el beneficio financiero
    del CVP1.
  \item \textbf{Deterioro de la experiencia del cliente:} la coexistencia de un canal digital
    inconsistente con un proceso de call center manual generaría contradicciones en el
    servicio, afectando directamente el NPS y acelerando la migración de clientes hacia los
    competidores.
  \item \textbf{Pérdida de la ventana competitiva del CVP2:} el segmento de 18 a 25 años que
    el CVP2 busca capturar no esperará. Competidores digitales ya operativos pueden ocupar ese
    espacio mientras Todo Tarjetas resuelve resistencias internas al cambio.
\end{itemize}

\subsubsection*{2. Mediano Plazo (1–3 años)}

\begin{itemize}
  \item \textbf{Deuda técnica en el asistente GenAI:} sin una cultura de mejora continua y
    despliegue iterativo, el modelo del CVP1 dejará de actualizarse y perderá precisión
    progresivamente, convirtiéndose en un sistema obsoleto que genera más fricción que la que
    resuelve.
  \item \textbf{Pérdida de participación en el segmento joven:} si el CVP2 no logra adopción
    en los primeros 18 meses por resistencia de los analistas de riesgo al scoring automático,
    los USD 1\,580\,000 en ahorros proyectados no se realizarán y el segmento 18–25 quedará
    consolidado en manos de competidores digitales.
  \item \textbf{Incapacidad para atraer talento especializado:} una organización percibida como
    culturalmente tradicional tiene dificultades estructurales para reclutar ingenieros de IA,
    diseñadores de experiencia digital y data scientists, perfiles críticos para sostener y
    evolucionar ambos CVPs en el tiempo.
  \item \textbf{Inversión tecnológica sin retorno:} los USD 2,14\,M comprometidos en
    plataforma, licencias de LLM y motor de originación digital generarán un retorno cercano
    a cero si la organización no adopta los comportamientos necesarios para operar y escalar
    los CVPs.
  \item \textbf{Riesgo de disrupción competitiva irreversible:} fintechs nativas digitales que
    ya operan en el Maturing Cluster pueden replicar las funcionalidades del CVP1 y CVP2 en
    menos tiempo del que le tomaría a Todo Tarjetas completar su transformación cultural,
    comprometiendo la relevancia de la compañía en el mercado de procesamiento digital.
\end{itemize}

\bigskip

La evolución cultural es la condición habilitante de ambos CVPs. El CVP1 requiere una cultura
ágil, tolerante al fallo y orientada a la mejora continua, pues sin ella el primer error del
asistente GenAI se convertirá en argumento para abandonar el proyecto antes de que demuestre
su valor real. Por su parte, el CVP2 requiere una cultura \textit{data-driven} que genere
confianza en el scoring automático y libere a los analistas de riesgo de la validación manual,
ya que sin ese cambio de mentalidad el proceso de originación digital seguirá siendo tan lento
y costoso como el modelo físico que busca reemplazar.

\end{document}
